The course notes classify musical instruments into string, percussion, and wind instruments (p. 88). In each case, a mechanical oscillator drives pressure fluctuations in the surrounding air. The oscillator and its boundary conditions determine the frequency content, while the instrument body or air column couples that motion to the listener.

\textbf{1. Methods of exciting a string.} A string can be plucked, bowed, or struck with a hammer. Plucking gives the string an initial displacement and releases it, bowing supplies repeated frictional excitation, and hammering supplies a short impact. In all three cases, the string vibrates in standing-wave modes and transfers energy through the bridge to a soundboard or instrument body, which radiates the sound more effectively into the air.

\[
\boxed{\text{Pluck, bow, or strike the string.}}\qquad\text{(A.6.1)}
\]

\textbf{2. Methods of exciting a wind instrument.} The two broad methods are edge-tone excitation and reed excitation. In an edge-tone instrument, a focused air jet alternates across a sharp opening and excites the air column, as in a flute. In a reed instrument, airflow makes a flexible reed or the player's lips open and close the passage periodically, as in a clarinet, oboe, or trumpet. The resulting pressure disturbance couples to resonant modes of the air column \cite{WolfeNDWoodwind,WolfeNDBrass}.

\[
\boxed{\begin{gathered}
\text{Edge-tone excitation by an air jet, and}\\
\text{reed or lip excitation by a periodically opening passage.}
\end{gathered}}\qquad\text{(A.6.2)}
\]

\textbf{3. String instruments excited by a hammer.} Examples include the piano and hammered dulcimer. The course notes describe the hammered dulcimer as a percussion-stringed instrument whose strings are struck by small hammers and whose resonant soundboard radiates the resulting vibration (pp. 170--171).

\[
\boxed{\text{Piano and hammered dulcimer.}}\qquad\text{(A.6.3)}
\]

\textbf{4. Major categories of musical instruments.} The three categories used in the course notes are string, percussion, and wind instruments.

\textbf{String instruments.} A stretched string is excited by plucking, bowing, or impact. The string supports transverse standing waves, and its natural frequencies depend on its length, tension, and mass per unit length. The bridge and soundboard couple the string motion to the air and strengthen the radiated sound.

\[
\boxed{\text{String: vibrating stretched strings and their standing-wave modes.}}\qquad\text{(A.6.4a)}
\]

\noindent\begin{minipage}{\linewidth}\color{black}
\textbf{Percussion instruments.} A membrane, bar, plate, or instrument body is struck, shaken, or scraped. The short interaction supplies an impulse that excites the object's natural modes. Its mass, stiffness, geometry, and boundary conditions determine the modal frequencies, while damping controls the decay of the sound.

\[
\boxed{\text{Percussion: impact or motion excites structural vibration modes.}}\qquad\text{(A.6.4b)}
\]
\end{minipage}\par

\textbf{Wind instruments.} A jet, reed, or vibrating lips excites an air column inside a tube. The pressure and particle-displacement waves reflect from the tube boundaries and form standing modes. Open and closed ends impose different pressure and displacement conditions, so the tube length, openings, valves, and bore determine the resonant frequencies. The course notes use the vibrating reed as a wind-instrument application of cantilever-beam vibration (p. 211).

\[
\boxed{\text{Wind: an air-column resonance driven by an edge jet, reed, or lips.}}\qquad\text{(A.6.4c)}
\]

\textbf{5. Common string instruments.} Common examples are the violin, viola, cello, double bass, guitar, harp, piano, and hammered dulcimer. The piano and hammered dulcimer are string instruments even though their strings are excited by hammers rather than by fingers or a bow.

\[
\boxed{\begin{gathered}
\text{Violin, viola, cello, double bass, guitar, harp,}\\
\text{piano, and hammered dulcimer.}
\end{gathered}}\qquad\text{(A.6.5)}
\]
